\begin{problem}[33届物理竞赛复赛]{59}
    蹦极是年轻人喜爱的运动。为研究蹦极过程，现将一长为 L 、质量为 m 、当仅受到绳本身重力时几乎不可伸长的均匀弹性绳的一端系在桥沿 b ，绳的另一端系一质量为 M 的小物块（模拟蹦极者）；假设 M 比 m 大很多，以至于均匀弹性绳受到绳本身重力和蹦极者的重力向下拉时会显著伸长，但仍在弹性限度内。在蹦极者从静止下落直至蹦极者到达最下端、但未向下拉紧绳之前的下落过程中，不考虑水平运动和可能的能量损失。重力加速度大小为 g 。

    \begin{subproblem}
        求蹦极者从静止下落距离 $y$ （ $y < L$ ）时的速度和加速度的大小，蹦极者在所考虑的下落过程中的速度和加速度大小的上限。
    \end{subproblem}
    \begin{subproblem}
        求蹦极者从静止下落距离 $y$ （ $y < L$ ）时，绳在其左端悬点 $\mathbf{b}$ 处张力的大小。
    \end{subproblem}
\end{problem}